\exercisesection{3}

\begin{enumerate}
\def\labelenumi{\arabic{enumi}.}
\item
  在第 1 节命题 1 的假设下，证明每个素理想 \(\mathfrak{p} \in \operatorname{Ass}_f(\mathbf{M})\) （§ 2，习题 17）都是分次的，并且是包含 M 的某个齐次元素的零化子的素理想集的极小元。（先证明最大的分次理想 \(\mathfrak{p}'\) （含于 A 的素理想 \(\mathfrak{p}\) 中的）是素的，注意：若两个元素 \(x, y\) 的乘积属于 \(\mathfrak{p}'\) ，则它们最高次数的齐次分量（ \(\neq 0\) ）的乘积也如此。再注意，若一个素理想包含某个元素 \(z \in \mathbf{M}\) 的零化子，则它也包含诸齐次分量（ \(\neq 0\) ）之属于 \(z\) 者中至少一个的零化子。）
\item
  在型为 \(\Delta\) 的分次环 A 中（其中 \(\Delta\) 无挠），设 p 是 A 的素理想集的一个极小元，它必然是分次的（习题 1）。证明对每个齐次元素 \(t \in p\) ，存在齐次元素 \(s \notin p\) 及整数 n \textgreater{} 0 使 \(st^{n} = 0\) 。由此把命题 3、4 和 5 推广到未必 Noether 的分次环（采用 §2 习题 12 与 20 中的定义）。
\item
  设 A 是型为 \(\Delta\) 的分次环（其中 \(\Delta\) 无挠），M 是强 Lasker 有限生成分次 A-模（ \(\S 2\) ，习题 28）。证明子模 \(Q(\mathfrak{p}) = Q(\mathfrak{p}, e_{\mathfrak{p}}(\mathbf{M}))\) （其中 p 遍历 \(\operatorname{Ass}_{f}(\mathbf{M})\) ）（ \(\S 2\) 习题 4 中定义的）都是分次子模。
\end{enumerate}

